\clearpage\SourceLeaf{199}%
\pdfbookmark[1]{Gauss to Schumacher, 23 November 1847}{schumacher18471123}
\fancyhead[L]{\small GAUSS TO SCHUMACHER}
\begin{center}
{\small GAUSS TO SCHUMACHER.\par}
\medskip\hrule\medskip
{\small Astronomische Nachrichten edited by H.\ C.\ \textsc{Schumacher}.\par
Vol.\ 27.\ No.\ 625.\ p.\ 1.\par}
\medskip\hrule
\end{center}\medskip
\begin{flushright}Göttingen, 23 November 1847.\end{flushright}
The second set of elements for Iris, printed in no.\ 615 of the A.\ N., was calculated by Prof.\ \textsc{Goldschmidt} just after 19 September. It continued to agree well with the observations in the following months (on 17 and 18 November the difference from Mr.\ \textsc{Rümcker}’s meridian observation was $16''$ in right ascension and $2''$ in declination). One is therefore justified in regarding the elements as already a close approximation. I accordingly asked Prof.\ \textsc{Goldschmidt} to use them to calculate the zodiac of Iris. I give the result of this work below, as it may make it easier to establish whether earlier observations were of the same object.

Taking another look on this occasion at the paper in which I gave the general method for determining the boundaries of such a zodiac 44 years ago, I see that I already indicated there the exceptional cases where the region on the celestial sphere in which a body orbiting the Sun according to \textsc{Kepler}’s laws can appear geocentrically has either no boundary at all or only \textit{one}, although in the first case the method gives a line that returns into itself and in the other two such lines. I also raised the question there of what, in such exceptional cases, these lines found by calculation actually
\SourceLeaf{200}%
mean. I confined myself to those hints at the time because going into more detail there would have been an hors d’oeuvre, and because I wanted to leave others the pleasure of working on a mathematical problem that I thought was not uninteresting. But I have heard of no one taking up the question in the long interval, during which it had quite slipped my mind. So I take this opportunity to set out at least the backbone of what is needed to answer it.

A geocentric position of the planet (or comet) in question is obtained by combining a point on its orbit with a point on the Earth’s orbit. But two, three, or four different combinations can also give the same geocentric position. To make this easier to picture, pass a plane through the Sun, perpendicular to the straight line that represents a given geocentric position on the celestial sphere. Project both the Earth’s orbit and the planet’s orbit orthographically onto this plane. Both projections are ellipses, or more generally conic sections, which will intersect either not at all, twice, or four times. A tangency is to be treated as the merging of two intersections. Each intersection represents a geocentric position that is either identical with the given one or opposite it on the celestial sphere, depending on whether the first or the second of the two points on the planet’s orbit and the Earth’s orbit whose projections coincide lies higher. Take the upper side of the projection plane to be the side on which the given geocentric position lies.

It is clear from this that if any point on the celestial sphere. is proposed as a possible geocentric position, there may be no way at all to obtain it as such a position, or there may be \textit{one}, two, three, or four ways. For a given planet, the whole surface of the celestial sphere thus falls into different regions. In general, the lines determined by the method I gave separate two regions of the sphere’s surface such that the points on one side can be geocentric positions in two more ways than those on the other side. The points on the dividing lines mark the transition. That is, they can be geocentric positions in \textit{one} more way than the points on one side and in \textit{one} fewer way than the
\SourceLeaf{201}%
points on the other side. It is also possible to give a criterion for deciding a priori on which side of the dividing line there are two more solutions than on the other side, though I will leave that aside for now.

The orbits of all planets known so far are such that the dividing lines found by the theory are always actual boundaries. There are two dividing lines that divide the celestial sphere into three regions. Two are isolated regions containing no geocentric positions at all. The third lies between them like a belt and contains all the geocentric positions. Every point inside the belt can be a geocentric position in two ways, and every point on its boundary in \textit{one} way.

But one can imagine hypothetical orbits (several comet orbits may be cases in point, though I have made no investigations so far) where two boundaries likewise divide the celestial sphere into three regions. One contains no geocentric positions at all, another contains the geocentric positions obtainable in two ways, and the third contains points that can be geocentric positions in four different ways.

Further variety arises when \textit{one} dividing line crosses itself once or twice, forming a single or a double loop. In the first case it cuts off two regions from the belt-shaped part, and in the second three, with respectively 4 and 0 or 4, 0, 4 ways of obtaining solutions.

An orbit that links with the Earth’s orbit like a link in a chain behaves differently again. In that case there is only \textit{one} connected boundary line and thus two separate regions of the surface. The points of one region can be geocentric positions in \textit{one} way, and those of the other in three ways.
\par\medskip\centerline{\rule{.25\textwidth}{.3pt}}
